% chapters/chapter3/3_6_suggestions.tex
% !TeX root=../../main.tex

% Suggestions for Future Work
\section{پیشنهادها برای تحقیقات آتی}

با توجه به دستاوردهای این پژوهش در توسعهٔ معماری کنترل یکپارچه برای سیستم‌های جرثقیل سقفی، مسیرهای زیر جهت گسترش و ارتقای این سامانه در مطالعات آینده پیشنهاد می‌شود:

\begin{itemize}
    \item \textbf{پیاده‌سازی آزمایشگاهی و ارزیابی تجربی:} اجرای الگوریتم \lr{NMPC} بر روی یک پروتوتایپ آزمایشگاهی به منظور بررسی پایداری سیستم در برابر نویز سنسورها، تأخیرهای ارتباطی در حلقهٔ فیدبک و عدم‌قطعیت‌های مدل‌سازی (نظیر اصطکاک‌های غیرخطی در مفاصل و مسیر ریل).
    \item \textbf{ارتقای مدل به آونگ دوگانه (\lr{Double Pendulum}):} در بسیاری از کاربردهای سنگین صنعتی (مانند انتقال کانتینرها)، توزیع جرم قلاب و بار باعث بروز دینامیک آونگ دوگانه می‌شود. بسط معادلات غیرخطی و تنظیم مجدد تابع هزینهٔ کنترل‌کننده برای مهار نوسانات در دو فرکانس طبیعی متفاوت، گامی حیاتی در صنعتی‌سازی این پژوهش است.
    \item \textbf{ادغام الگوریتم‌های پرهیز از مانع (\lr{Obstacle Avoidance}):} بهره‌گیری از سامانه‌های بینایی ماشین برای تشخیص موانع متحرک در محیط کارگاه و تزریق مستقیم این داده‌ها به عنوان قیود متغیر با زمان در مسئلهٔ بهینه‌سازی \lr{NMPC}، تا سیستم بتواند مسیر ایمن را به صورت برخط بازطراحی کند.
    \item \textbf{کاهش بار محاسباتی با رویکردهای یادگیری ماشین:} استفاده از شبکه‌های عصبی عمیق برای تقریب قانون کنترل صریح (\lr{Explicit MPC}) یا به‌کارگیری الگوریتم‌های یادگیری تقویتی (\lr{Reinforcement Learning}) به منظور استنتاج سریع سیگنال‌های کنترلی، که امکان استقرار این معماری را بر روی سخت‌افزارهای صنعتی با توان پردازشی محدود فراهم می‌سازد.
\end{itemize}